\setcounter{section}{7}

\section{Elektronen im Festkörper - Energiebänder}
\subsection{Adiabatische Näherung}
Die Bewegung der Rumpfionen und der Elektronen im Festkörper können unabhängig voneinander betrachtet werden. 
Im Rahmen der mean-field approximation betrachten wir näherungsweise ein einzelnes Elektron im effektiven Potential aller anderen Teilchen.
Zur Beschreibung reicht hier also ein Ein-Elektron-Hamiltonoperator.

Diesen setzen wir in die Schrödingergleichung ein und erhalten
\stepcounter{equation}
\begin{equation}
	\label{eq:II.8.2}
	\left[ - \frac{\hbar^2}{2m} \Delta + V(\Vec{r}) \right] \psi(\Vec{r}) = E \psi(\Vec{r})
\end{equation}
mit dem Gitterperiodischen Potential
\begin{equation}
	\label{eq:II.8.3}
	V(\Vec{r}) = V(\Vec{r} + \Vec{t}),
\end{equation}
wobei die geneaure Form dieses Potentials vorerst nicht von Interesse für uns ist. $\vec t$ ist dabei ein beliebiger Translationsvektor im Gitter mit
$$
\vec t = t_1 \Vec{a}_{1} + t_2 \Vec{a}_{2} + t_3 \Vec{a}_{3}.
$$ 
Da das Potential $V(\Vec{r})$ gitterperiodisch ist, lässt es sich in eine Fourierreihe entwickeln mit
\begin{equation}
	\label{eq:II.8.4}
	V(\Vec{r}) = \sum_{\Vec{G}}^{} V_{\Vec{G}} e^{i \Vec{G} \Vec{r}},
\end{equation}
wobei $\Vec{G} = h\Vec{g}_{1} + h\Vec{g}_{2} + l\Vec{g}_{3}$ ein reziproker Gittervektor ist. 
$V_{\Vec{G}}$ ist ein für den Kristall charakteristischer Fourierkoeffizient.
 $\psi(\Vec{r})$ kann als Entwicklung nach ebenen Wellen angesehen werden mit
\begin{equation}
	\label{eq:II.8.5}
	\psi(\Vec{r}) = \sum_{\Vec{k} }^{} C_{\vec k} e^{i \Vec{k} \Vec{r}} 
\end{equation}

Die periodischen Randbedingungen im Kristall erzwingen diskrete Werte für $k_{i}$, da \\${\psi_{\Vec{k}}(\Vec{r}) = \psi_{\Vec{k}}(\Vec{r} + \Vec{L})}$ mit der makroskopischen Länge des Kristallverbandes $L = n_{i} a$. 
Dabei ist $N$ die Zahl der Elementarzellen und $a$ die Gitterkonstanten.
Erlaubte Wellenzahlen sind somit
$$
\begin{aligned}
	k_{x} &= 0, \pm \frac{2\pi}{L}, \pm \frac{4\pi}{L}, \ldots&, \frac{2\pi n_{x}}{L} \\
	k_{y} &= 0,\ldots, &, \frac{2\pi n_{y}}{L} \\
	k_{z} &= 0, \ldots, &, \frac{2\pi n_{z}}{L}
\end{aligned}
$$ 
mit $n_{x}, n_{y}, n_{z} = 0, \pm 1, \ldots$. 

Einsetzen von \eqref{eq:II.8.4} und \eqref{eq:II.8.5} in \eqref{eq:II.8.2} liefert
$$
\sum_{\Vec{k} }^{} \frac{\hbar ^2 k^2}{2m} C_{\Vec{k}} e^{i \Vec{k} \Vec{r}} + \underbrace{\sum_{\Vec{k} \Vec{G}}^{} C_{\Vec{k}} V_{\Vec{G}} e^{i (\Vec{k} + \Vec{G}) \Vec{r}}}_{(\ast)}= E \sum_{\Vec{k}}^{}  C_{\Vec{k}} e^{i\Vec{k}\Vec{r}}.
$$ 
Durch die Umbennenung des Summationsindex $\Vec{k} \to \Vec{k} -\Vec{G}$ schreiben wir
$$
(\ast) = \sum_{\Vec{k} \Vec{G}}^{} C_{\Vec{k} - \Vec{G}} V_{\Vec{G}} e^{i\Vec{k} \Vec{r}}.
$$ 
Damit folgt
\begin{equation}
	\label{eq:II.8.6}
	\sum_{\Vec{k} }^{} e^{i\Vec{k} \Vec{r}} \left[ \left( \frac{\hbar ^2 k^2}{2m} - E \right) C_{\Vec{k}} + \sum_{\Vec{G} }^{}  V_{\Vec{G}} C_{\Vec{k} -\Vec{G}}  \right] = 0.
\end{equation}
Da es sich bei $\exp(i\Vec{k}\Vec{r})$ um einen orthogonalen Satz von Funktionen handelt, müssen alle Entwicklungskoeffizienten Null sein:

\begin{important}
	$$
	\left( \frac{\hbar ^2 k^2}{2m} - E \right) C_{\Vec{k}} + \sum_{\Vec{G} }^{} V_{\Vec{G}} C_{\Vec{k} -\Vec{G}} = 0  
	$$ 
\end{important}

Dieser Satz von algebraischen Gleichungen ist eine andere Darstelung der Schrödingergleichung im Wellenzahlraum ($\Vec{k}$-Raum).
Er liefert für jedes $\Vec{k}$ eine Lösung $\psi_{\Vec{k}}$ mit zugehörigem Eigenwert $E_{\Vec{k}}$.

Die Gleichungen koppeln nur solche Entwicklungskoeffizienten $C_{\Vec{k}}$, deren $\Vec{k}$-Vektoren sich um jeweils reziproke Gittervektoren $\Vec{G}$ von diesen unterscheiden:
$$
C_{k} \quad\leftrightarrow\quad C_{\Vec{k} - \vec G}\;,\;C_{\Vec{k}- \Vec{G}'}\;,\;\ldots
$$ 
Der Gleichungssatz enthält also Bestimmungsgleichungen von $C_{\vec k}$.

Jedes der $N$ Gleichungssysteme liefert eine Lösung, die sich als Superposition von ebenen Wellen darstellen lässt.
Zu einer Energie
$$E_{\Vec{k}} = E(k)$$
gehört eine Wellenfunktion
\begin{equation}
	\label{eq:II.8.8}
	\psi_{\Vec{k}}(\Vec{r}) = \sum_{G }^{} C_{\Vec{k} -\Vec{G}} e^{i(\Vec{k}-\Vec{G}) \Vec{r}}.
\end{equation} 
Die Wellenfunktion hat auch eine Formulierung nach dem
\begin{important}
	\textbf{Bloch'schen Theorem}
	\vspace{1ex}
	\begin{equation}
		\label{eq:II.8.9}
		\psi_{\Vec{k}} (\Vec{r}) = \sum_{\Vec{G}}^{} C_{\Vec{k} -\Vec{G}} e^{-i \Vec{G}\Vec{r}} e^{i\Vec{k}\Vec{r}} = u_{\Vec{k}} (\Vec{r}) e^{i\Vec{k}\Vec{r}} 
	\end{equation}
	mit 
	\begin{equation}
		\label{eq:II.8.10}
		u_{\Vec{k}}(\Vec{r}) = u_{\Vec{k}}(\Vec{r} + \Vec{t}) .
	\end{equation}
\end{important}

Die Lösung ist das Produkt aus einer ebenen Wellen und einer gitterperiodischen Funktion.

Die Lösung ist translationsinvariant:
$$
u_{\Vec{k}}(\Vec{r} + \Vec{t}) = \sum_{}^{} C_{\Vec{k}-\Vec{G}} e^{-i\Vec{G}(\Vec{r}+\Vec{t})} = \underbrace{e^{i\Vec{G}\Vec{t}}}_{=1} u_{\Vec{k}}(\Vec{r}) 
$$
$$
\Rightarrow U_{\Vec{k}}(\Vec{r}+\Vec{t}) = U_{\Vec{k}}(\Vec{r})
$$ 

\paragraph{Weitere Konsequenzen}\,\\ 
Die Wellenfunktion erfüllt insbesondere
\begin{equation}
	\label{eq:II.8.11}
	\psi_{\Vec{k}}(\Vec{r}) = \psi_{\Vec{k}+\Vec{G}} (\Vec{r}).
\end{equation}
Einsetzen in die Schrödingergleichung $H\psi_{\Vec{k}}(\Vec{r}) = E(\Vec{k}) \psi_{\Vec{k}}(\Vec{r})$ gibt
\begin{align}
	\nonumber
	\Leftrightarrow&& H \psi_{\Vec{k}+\Vec{G}} (\Vec{r}) =& E(\Vec{k}+\Vec{G}) \psi_{\Vec{k}+\Vec{G}}(\Vec{r})\\
	\nonumber
	\Leftrightarrow&&  H \psi_{\Vec{k}}(\Vec{r}) =& E(\Vec{k}+\Vec{G}) \psi_{\Vec{k}}(\Vec{r})\\
	\label{eq:II.8.12}
	\Leftrightarrow&&  E(\Vec{k}) =& E(\Vec{k} + \Vec{G})
\end{align}
Die Wellenfunktionen der Elektronen und ihre Eigenwerte wiederholen sich im $\Vec{k}$-Raum periodisch. 
Es reicht deshalb, sich auf Lösungen in der ersten Brillouinzone zu beschränken.

Es gibt zu jedem $\Vec{k}$-Vektor unendlich viele Lösungen $E_{n}(\Vec{k}) = E(\Vec{k}+\Vec{G}_{n})$, wobei $\Vec{G}_{n}$ der zum Energieeigenwert $E_{n}$ gehörige reziproke Gittervektor ist.\\
$n$ ist also der Bandindex mit
$$
E_{1}(\Vec{k}) \le E_2(\Vec{k}) \ldots E_{n}(\Vec{k}).
$$ 

\subsection{Bänderschema 1}
\subsubsection{Näherung von quasi-freien Elektronen}
Die \textbf{Bragg-Reflexion} ist ein charakteristisches Kennzeichen für die Ausbreitung von Wellen in Kristallen.
Deshalb ist sie die \textbf{Ursache der Energielücke}.


Wir betrachten einen kubischen Kristall mit der Netzebene $[100]$ und der Gitterkonstante $a$ (vgl. \autoref{fig:vl27_abbildung_1}).\\[1ex]
\begin{minipage}{0.5\linewidth}
	Wir erhalten bei einer \SI{180}{\degree}-Streuung aus der Bragg-Bedingung
	$$
	\begin{aligned}
		2d \sin\left( \frac{\theta}{2} \right) &= m \lambda \\
		2a &= m \lambda \\
		\lambda &= \frac{2a}{m} \\
		k = \frac{2\pi}{\lambda} &= \pm \frac{\pi}{a} m .\\
	\end{aligned}
	$$ 
\end{minipage}
\begin{minipage}{0.5\linewidth}
	\begin{figure}[H]
		\centering
		\def\svgwidth{.386\columnwidth}
		\import{figures/}{vl27_abbildung_1.pdf_tex}
		\caption{Netzebenen $[100]$ im Abstand $d$. In blau ist eine Streuung um \SI{180}{\degree} eingezeichnet.}
		\label{fig:vl27_abbildung_1}
	\end{figure}
\end{minipage}

Bei $k=\pm \pi /a$ gibt es keine laufenden Wellen, d.h. gleiche Anteile von nach links und nach rechts laufenden Wellen müssen bestehen. Es handelt sich also um eine stehende Welle.

Es existieren zwei Lösungen stehender Wellen mit
$$
\begin{aligned}
	\psi^{(+)} &= e^{i \frac{\pi}{a} x} + e^{- i \frac{\pi}{a} x} = \cos\left( \frac{\pi}{a} x \right),  \\
	\psi^{(-)} &= e^{i\frac{\pi}{a} x} - e^{-i \frac{\pi}{a} x} = 2i \sin \left( \frac{\pi}{a}x \right).  
\end{aligned}
$$ 

\begin{figure}[H]
    \centering
    \incfig{vl27_abbildung_2}
	\caption{Die Lösungen für stehende Wellen $|\psi^{(+)}|^2$ und $|\psi^{(-)}|^2$. Ihre physikalische Interpretation sind räumlich modulierte Ladungsdichten.}
    \label{fig:vl27_abbildung_2}
\end{figure}

\begin{figure}[H]
    \centering
    \incfig{vl27_abbildung_3}
	\caption{Energie-Dispersionsrelation im Band. 
		In grün sind die stehenden Wellenlösungen eingezeichnet.
		Wegen der Periodizität lassen sich alle stehenden Wellen in die erste Brillouinzone übertragen (rot).
		Zwischen den Bändern exisiteren Energielücken.
	}
    \label{fig:vl27_abbildung_3}
\end{figure}

\subsubsection{Bänderschema 2: Näherung fest gebundener Elektronen}
Idee: Stark an Atomrümpfe gebundene Elektronen werden ihre Eigenschaften auch im Kristallverband weitgehend behalten.

Für \textbf{ein Atom} ist die stationäre Schrödingergleichung
$$
\hat{H}_{A}(\Vec{r} -\Vec{r}_{n}) \varphi_{i} (\Vec{r}-\Vec{r}_{n}) = E_{i} \varphi_{i}(\Vec{r}-\Vec{r}_{n}) 
$$ 
mit dem Hamilton-Operator $\hat{H}_{A}(\Vec{r}- \Vec{r}_{n})$ für ein freies Atom am Gitterplatz $n$ und der Wellenfunktion $\varphi_{i}(\Vec{r}-\Vec{r}_{n})$ eines Elektrons im Zustand $i$ am Atom $r_{n}$.

Im \textbf{Kristallverband} betrachten wir die Potentiale der übrigen Atome als Störung $V$ zum Einatom-Hamiltonoperator $\hat{H}_{A}$:
$$
\hat{H} = \hat{H}_{A} + V(r-r_{n)} = - \frac{\hbar^2}{2m} \Delta + V_{A} (r-r_{n}) + V(r-r_{n})
$$ 
